% !TeX root = ../se200_identity_regimes.tex

% =============================================================================
\section{Formal Setting}
\label{sec:formal_setting}
% =============================================================================

The results in this paper are conditional on the assumptions fixed here.
These assumptions specify the operating conditions under which the
neutral-substrate constraint of Section~\ref{sec:background} is applied.
The goal is to identify the lower-bound identity structure required of a
neutral accountability substrate under those conditions.

% -----------------------------------------------------------------------------
\subsection{Persistent Disagreement}
\label{subsec:persistent_disagreement}
% -----------------------------------------------------------------------------

We assume environments in which disagreement is persistent and extensions are uncoordinated.
Extensions may introduce
incompatible causal accounts, incompatible normative attributions, and
incompatible interpretive classifications.
They cannot be forced to converge.
Substrate revision may not be used to resolve interpretive conflict.
Consequently, substrate-level admissibility must not depend on agreement among extensions.
Neutrality must be maintained without recourse to convergence.

Persistent disagreement is not mere temporary uncertainty.
It is the condition in which multiple admissible frameworks
may remain simultaneously usable while disagreeing about
causation, responsibility, justification, applicability, or normative significance.
A neutral substrate must therefore support shared reference
without embedding any one framework's interpretation as the substrate-level account.

% -----------------------------------------------------------------------------
\subsection{Stable Reference in Strong Form}
\label{subsec:stable_reference_strong}
% -----------------------------------------------------------------------------

Extension-stable reference is inherited from \emph{Neutral Substrates}.
Here we fix the demanding case in which stability must
survive ordinary transformations of records.
Identity conditions must not depend on interpretive convergence.
An entity that must change its ontological status to enter a
new relationship has not been stably referred to.
A stability-critical distinction may not be carried only by a
role, flag, or contextual predicate.
Such devices rely on downstream conventions and can therefore
permit divergent judgments under persistent disagreement.

This is a deliberately strong setting.
Engineering systems often manage identity pressure through
negotiated, local, or context-relative schemes, which may be useful in
practice but are not the case analyzed here.
This paper studies the case in which stable reference must be available
as common ground before interpretive agreement is available.

% -----------------------------------------------------------------------------
\subsection{Minimality}
\label{subsec:minimality}
% -----------------------------------------------------------------------------

We minimize referential structure subject to the above constraints.
Minimality is measured by the number of identity regimes
needed to preserve neutrality and stable reference.
A reduction is admissible only when it preserves the relevant identity distinctions
without relocating them into hidden discriminators.

A reduction counts as genuine only when it does not reintroduce
a collapsed distinction through
a role, flag, contextual predicate, workflow state, schema convention,
or application-level discriminator.
The criterion is analytic.
It identifies the regime structure required under the stated constraints.
It does not recommend use of the core set beyond
the neutral-substrate conditions analyzed here.

% -----------------------------------------------------------------------------
\subsection{Transformation Basis}
\label{subsec:transformation_basis}
% -----------------------------------------------------------------------------

Identity pressure appears when records are transformed.
We therefore fix a finite transformation basis
\[
\Fset =
\{\mathrm{RE},\mathrm{AN},\mathrm{RF},\mathrm{AD},\mathrm{RC},
\mathrm{RA},\mathrm{SU},\mathrm{BF},\mathrm{PV},\mathrm{SE}\}.
\]
The basis is not claimed to be complete for all possible record systems.
It is a representative basis of ordinary operations in
provenance, versioning, institutional recordkeeping, and data-governance practice.
The lower-bound result is conditional on this basis.
Adding further transformations cannot collapse distinctions
already witnessed by this basis.
It may leave the lower bound unchanged or force additional splits.

\begin{definition}[Transformation basis]
\label{def:transformation_basis}
The transformation basis $\Fset$ consists of the following admissible operation families:
\begin{enumerate}[label=(\alph*),nosep]
  \item $\mathrm{RE}$, re-expression: change of representation format while preserving the represented content;
  \item $\mathrm{AN}$, annotation: addition of non-identity-bearing information;
  \item $\mathrm{RF}$, refinement: increase in descriptive, structural, or normative detail;
  \item $\mathrm{AD}$, aggregation or decomposition: restructuring into coarser or finer components;
  \item $\mathrm{RC}$, re-contextualization: change in applicability context;
  \item $\mathrm{RA}$, reassignment: change in association, bearer, or assignment relation;
  \item $\mathrm{SU}$, substitution: replacement of the identity carrier;
  \item $\mathrm{BF}$, branch or fork: divergence into multiple continuations;
  \item $\mathrm{PV}$, provenance extension: addition of historical, derivational, or evidentiary metadata;
  \item $\mathrm{SE}$, state evolution: change in state of an enduring referent over time.
\end{enumerate}
\end{definition}

The transformation families are structural rather than interpretive.
They do not assert causal or normative claims.
They describe operations that may be logged or recognized without
deciding what those operations mean under any particular admissible framework.
Their significance for this paper lies only in their effect on identity.

% -----------------------------------------------------------------------------
\subsection{Classification of Transformations}
\label{subsec:classification_transformations}
% -----------------------------------------------------------------------------

A regime classifies each applicable transformation by how it affects the declared identity basis.
We use three classifications.

\begin{definition}[Transformation classification]
\label{def:transformation_classification}
For a regime $\tau$ and transformation family $f\in\Fset$, the classification
\[
\Class_\tau(f)\in\{\PRS,\BRK,\NEU,\NA\}
\]
has the following meaning:
\begin{enumerate}[label=(\alph*),nosep]
  \item $\PRS$: the transformation is identity-preserving for $\tau$;
  \item $\BRK$: the transformation is identity-breaking for $\tau$;
  \item $\NEU$: the transformation is identity-neutral for $\tau$;
  \item $\NA$: the transformation is inapplicable to $\tau$.
\end{enumerate}
\end{definition}

The distinction between $\PRS$ and $\NEU$ matters.
A $\PRS$ transformation preserves identity in a way that is relevant to the regime's identity basis.
An $\NEU$ transformation does not bear on that identity basis.
Both are non-breaking.
Only $\BRK$ breaks identity.
Thus $\NEU$ does not mean that the operation is ignored by the system.
It means that the operation does not alter identity under the regime being considered.
A fourth value, $\NA$, marks a transformation that cannot arise for the regime's carrier at all,
as distinct from one that arises but leaves identity untouched.

\begin{definition}[Non-breaking transformation]
\label{def:nonbreaking}
A transformation family $f$ is \emph{non-breaking} for regime $\tau$ when
\[
\Class_\tau(f)\in\{\PRS,\NEU\}.
\]
It is \emph{breaking} for $\tau$ when
\[
\Class_\tau(f)=\BRK.
\]
\end{definition}

% -----------------------------------------------------------------------------
\subsection{Identity Regimes}
\label{subsec:identity_regimes}
% -----------------------------------------------------------------------------

The new primitive of this paper is the identity regime.
An identity regime is not an ontological kind.
It is a transformation profile over an identity basis.

\begin{definition}[Identity regime]
\label{def:identity_regime}
An \emph{identity regime} $\tau$ consists of an identity basis
together with a classification function
\[
\Class_\tau:\Fset\to\{\PRS,\BRK,\NEU,\NA\}.
\]
The identity basis states what counts as sameness for the referent.
The classification function states how each transformation family affects sameness under that basis.
\end{definition}

\begin{definition}[Induced identity relation]
\label{def:induced_identity}
Let $\tau$ be an identity regime.
The induced identity relation $\sim_\tau$ is the equivalence relation
generated by transformations that are non-breaking for $\tau$.
Thus $r\sim_\tau s$ means that $r$ and $s$ are representations of the same referent under regime $\tau$.
\end{definition}

This definition makes identity transformation-relative.
The same pair of representations may coincide under one regime
and diverge under another.
The substrate must declare which identity basis is being preserved.

% -----------------------------------------------------------------------------
\subsection{Hidden Regimes}
\label{subsec:hidden_regimes}
% -----------------------------------------------------------------------------

A hidden regime is a collapsed identity distinction recovered outside
the declared substrate structure.
Hidden regimes obstruct genuine minimality.
They allow a system to appear to use fewer identity regimes
while relying on downstream discriminators to perform the missing identity work.

\begin{definition}[Hidden regime]
\label{def:hidden_regime}
A \emph{hidden regime} is a stability-critical distinction that
is not represented as an identity regime in the substrate,
but is instead recovered through a role, flag, contextual predicate,
workflow state, schema convention, or application-level discriminator.
\end{definition}

Hidden regimes violate the neutral-substrate discipline because their
application is not fixed by the substrate.
They require agreement about how the discriminator should be interpreted.
Under persistent disagreement, that agreement cannot be assumed.
A collapse that can be repaired only by a hidden regime is therefore
not a genuine reduction in declared referential structure.

% -----------------------------------------------------------------------------
\subsection{Scope of the Lower Bound}
\label{subsec:lower_bound_scope}
% -----------------------------------------------------------------------------

The lower bound proved below is conditional.
It holds for neutral accountability substrates satisfying the assumptions
of this section and using the transformation basis of Definition~\ref{def:transformation_basis}.
It asserts that, under the stated constraints,
fewer than nine regimes cannot preserve stable reference
without changing transformation behavior or reintroducing a hidden regime.
The claim is a lower bound only:
it does not assert that nine regimes
are sufficient for every substrate,
nor that the nine derived below exhaust
the coherent identity regimes.

The next section states the reference inventory.
The inventory supplies the referents that must remain addressable.
The transformation basis supplied here determines
which of those referents can share an identity regime.
